\section{Scope, theorem inputs, and the central reductions}

The main unrestricted tasks remain open in this manuscript: an explicit polynomial-size hitting set for general width-$w$ ROABPs, a uniformly efficient positive rule for all rectangular and plethysm multiplicities, and a determinant-versus-padded-permanent separation in the required asymptotic regime. Counterexamples below refute specified transfers, not those open problems. Every stated deduction has an identifier. The accompanying event register distinguishes complete proof readings, finite exact computations, dependency checks, and carried background. No Lean reconstruction or formal verification of the new mathematics is claimed.

The principal positive results are exact rather than heuristic. R-01--R-09 identify the CUE character and ratios operations, positive projector formulas, orbit-closure Gram ranks, and the missing quotient numerator. R-10 closes the height-two rectangle and quadratic determinant case. G-05 excludes every at-most-three-row multiplicity obstruction. G-07/R-12 give an exact four-row permanent-minus-determinant shifted-plethysm formula at size four. R-13 and the uniform endpoint repair R-27 now determine the quaternary permanent-image dimension in every size:
\[
 c_m=\min\left\{\binom{m+3}{3},4m^2-2m+2\right\}.
\]
R-24 supplies the exact size-seventeen base certificate; R-25 gives the border induction; R-26 supplies the uniform ordinary-character construction; R-27 recovers its six missing endpoint directions by a checked finite jet certificate and exact central insertion. The single graded Tor defect of R-23 consequently vanishes for every $m\ge17$. G-21 extends maximal parameter dimension to every $r\ge4$: $\dim\mathcal P_{m,r}=rm^2-2m+2$ for $m\ge17$, with the specified low-degree Tor vanishing.

G-17 classifies the strict four-row directions by $\nu_4>0$ and $\nu_1>(n-m)s$, and provides explicit first-stretch and progression bounds. G-18 establishes occurrence on both orbit closures and the all-integer asymptotic ratio. G-20 now gives its exact eventual sign for every pair $n>m\ge3$. For $m\ge17$, permanent dominance holds precisely when $n^2\le2m^2-m+1$, and determinant dominance holds when $n^2\ge2m^2-m+2$, in each fixed strict direction. There is no unresolved integer dimension-equality case. The exact four-row rectangular padding range is $m<n<4m/3$ for every $m\ge14$, and fourteen is sharp for this uniform eventual statement. This does not exclude individual earlier-degree gaps on the opposite side. G-16 specifies a nonrectangular weight family throughout the positive dimension range. G-15 proves dominant third jets and the necessary first-row condition $\lambda_1\ge d(n-m)+4$, while excluding full fourth-jet dominance. R-21 gives an exact orbit-multiplicity algorithm for fixed ambient dimension, with numerical-degree complexity. G-19 identifies one rectangular discriminant copy with zero determinant energy and positive permanent energy in equal size, including an explicit smooth cubic slice and five-row minimality. That specified-copy equation does not compare total multiplicities and vanishes under positive padding.

The higher-row construction G-22 uses an ordered Fourier frame and
cocharacter minimum weights to provide a common stable tensor in
every strict direction with $r\le m$. R-28 removes its scalar-orbit
period and gives $\dim\mathcal D_{n,r}=(r-2)n^2+2$.
For $m\ge17$ and $4\le r\le2m$, G-23 and its R-29/G-25 extension consequently prove
\[
 \frac{m_{j\nu}(\ell^{n-m}\operatorname{per}_m)}
      {m_{j\nu}(\det_n)}
 =\kappa_\nu j^{rm^2-(r-2)n^2-2m+r-1}(1+O(j^{-1})).
\]
Both multiplicities are positive at every sufficiently large integer
stretch, and the exponent is never zero. Its sign, with explicit
bounds for both alternatives, gives the exact eventual comparison
in every stated strict direction. G-24 completes the rectangular boundary comparison: the exact eventual positive range is $m<n<rm/(r-1)$. At equality the determinant eventually dominates, and above it every permanent rectangle vanishes. For the added rows $r>m$, the only positive parameter case is $r=n=m+1$: the explicit direction $(3,2^{m-1},1)$ has ratio $\kappa j(1+O(j^{-1}))$ in degree $2j$. G-26 additionally settles every nonrectangular first-row equality face through $m$ rows: Pieri bounds its multiplicity by the tail partition in one fewer variable, and the determinant eventually dominates. Below that face the padded multiplicity is zero. Thus every fixed direction through $m$ rows is classified eventually. R-32/R-33 and G-27/G-28 additionally classify every fixed nonrectangular direction through $2m$ rows. The lower comparison algebra may have a degree period and need only satisfy a quotient-dimension upper bound. Boundary comparison extends further to $r\le\min\{2n-2,m^2\}$ for nonrectangular labels. Rectangular directions beyond $m$ rows, larger strict directions and individual finite degrees retain their separate scope. R-31 proves a scalar-period obstruction for the unchanged two-product family at $2m+1$ rows in specified strict directions. Its odd-degree vanishing is a subcone statement, not a statement about either full orbit closure.


On isolation, I-01 gives the genuine Hoza transfer for independent vertex variables. I-03--I-04 identify coefficient flags and projected matroid cones; I-14 gives the weaker joint-span condition. The rank-profile construction I-21 and endpoint refinement I-30 supply an explicit six-pair interpolation hitting generator for width at most three in one unknown variable order, in every characteristic with the stated field-node convention. I-29 treats all words having one fixed singular interval rank. The full-matrix cone proof I-31 gives polynomial-size hitting sets for every fixed width and a pair budget controlled by the number of distinct interval ranks. Complete proofs are given here; the rank-profile and cone mechanisms are credited to a recent unrefereed preprint, whose broader affine-orbit claims are not assumed. I-32 refutes a polynomial coordinate-cone deletion bound even at a common center with every interval rank preserved. I-33 instead gives explicit small prime-cocharacter families for two dense products, their exact limitation, and a scalar product-closure repair. I-34 extends the prime family to bounded local nullity among separated products, with an exact three-versus-four-term cylinder distinction. P-14 uses scalar bridges to give $w-1$ interpolation pairs for factors of rank one or full rank, including one pair at width two. The unresolved target is a bound polynomial jointly in the growing width, length and degree; even a polynomial pair budget does not by itself meet that sampling target. I-35 weakens bridge preservation to equality of rank-one annihilator base loci, proves independent-seed composition, and isolates the missing synchronization condition with explicit counterexamples. P-16 gives an explicit degree16, width-exactly-three obstruction to the uncentered two-pair map on eighteen specified algebraic nodes. Its transition factors are generically invertible, and P-14 proves that the centered two-pair map hits the same polynomial. I-36 proves a pairwise-dominant maximal-minor preservation lemma and improves the bound to five pairs when every variable-dependent width-three factor is singular. I-37 reduces this to two centered pairs for the affine singular-factor subclass, in every characteristic and unknown order, by an irreducible-factor reduction and a tridiagonal second-coefficient argument. I-38 extends the same two-pair conclusion to affine words with isolated invertible factors: no two consecutive variable-dependent factors have rank three. Longer invertible runs and nonlinear singular factors remain open for this two-pair argument. I-39 refutes one-pole rank preservation even for an irreducible mixed affine scalar, and separately refutes the proposed abstract Hessian criterion while retaining its integrability limitation. P-17 gives characteristic-free coefficient-cut obstructions, including every balanced two-pair cubic annihilator through ten variables. I-40 bounds full-width invertible runs of a homogeneous translation-invariant tensor by its integer weight $n-2D$. At weight one this gives isolated invertible factors and closes the complete balanced two-pair width-three problem through nine variables, with an explicit finite rational sampler. I-41 uses the saturated weight-two support spaces and their actual divided-square dual pairing to preserve the four-coordinate connectors. Its bandwidth-three argument closes the balanced theorem through ten variables, with the same finite grid. The ordinary unweighted centered map is not identified with that balanced map for every ten-variable polynomial. P-18 proves the trilinear class and a smallest-counterexample support reduction. P-19 closes both remaining boundary cases: every multilinear width-three cubic is hit by balanced two-pair evaluation, in every characteristic and with no restriction on the number of variables, using $20(n+1)^2$ samples. Any remaining balanced width-three annihilator has its homogeneous degrees in the interval $[4,n-4]$. I-42 uses the actual rank-one reset structure to remove the Hessian step through four invertible factors in every characteristic, and five outside characteristic two; it supplies positive saturated-run classes. I-43 closes every weight-three connector: a Grassmann orbit line and the actual adjacent support bound force a rank-one update pattern, while the weight-zero anchor case is handled directly. Shared-parameter contractions then give the complete balanced two-pair theorem through eleven variables in every characteristic. I-44 classifies every chain anchor and preserves all interval ranks through five full-rank cores. The reset criterion yields weights at most four in every characteristic and at most five away from two; the complete balanced theorem follows through twelve variables, and thirteen away from two. Weight five in characteristic two and the next higher-degree cases remain open. The precise conditional matrix-invariant consequence remains I-08.

Purity gives sparse constructions, an actual-sheaf Gram criterion with full conductor accounting, and limitations of energy approximation. P-11 identifies explicit common arithmetic evaluation points with hitting-set encodings in its rational model. P-12 constructs such a point for the width-three class, together with an integer hitting curve and an exact trace-dual coefficient-span map. P-13 states the general centered-interpolation encoding and reduces the width-three field degree to $(nD+1)^{13}$ using I-30. It applies directly to the rank-spectrum construction I-31 with degree $(nD+1)^{2K_{\mathcal R}+1}$. These constructions use elementary algebraic number theory and require no additional RH input.

\subsection{Stipulated analytic inputs and supplied material}

We accept Alp\"oge's RH and Dirichlet GRH as theorem inputs, exactly as requested. In \texttt{ALPOGE\_Main.lean}, the exported declarations are \texttt{riemannHypothesis}, \texttt{riemannZeta\_eq\_zero\_re}, \texttt{LFunction\_ne\_zero\_of\_half\_lt\_re}, \texttt{LFunction\_eq\_zero\_re}, and \texttt{LFunction\_eq\_zero\_re\_of\_isPrimitive}; the predicates are defined in \texttt{ALPOGE\_Statements.lean}. The attached note supplies narrative descriptions of its numbered steps. Neither the imported Lean dependency tree nor the claimed all-files scan is in this attachment, and neither was rebuilt in this run. This is an input convention, not a review of the theorem's truth.

Masullo's general number-field Hecke GRH is used only conditionally, as the prompt's attribution to section~1.2. The supplied statement index has no numbered section~1.2 theorem or its proof. Accordingly we give no invented locator. Its Theorem~13.3.1.1, indexed PDF page~522, is a trace-formula claim; Conjecture~14.1.2.1, indexed page~529, concerns automorphic $\mathrm{GL}_n$ GRH and simplicity and remains explicitly conjectural there. These are different statements. The zero-statistics section names the precise hypotheses discharged by the stipulated input and those that remain.

The archive was unpacked once. The source manifest records the actual supplied files and hashes. The five cote~133 files are LaTeX transcriptions, not original page images. Assertions in the accompanying earlier study about visual readings remain credited to that study; our transcription table records mathematical necessity separately from paleographic confidence.

\subsection{Why the three multiplicity problems are distinct}

Let $V=\mathbb C^{n^2}$, $W=\operatorname{Sym}^nV^*$, and $A=\operatorname{Sym}(W^*)$. For a homogeneous form $f\in W$, write
\[
 m_f(\lambda,d)=\operatorname{mult}_{S_\lambda V}\mathbb C[\overline{GL(V)f}]_d,
 \qquad a_\lambda(d[n])=\operatorname{mult}_{S_\lambda V}A_d.
\]
The symbol $g(\lambda,\mu,\nu)$ denotes the ordinary Kronecker coefficient. One must distinguish the ambient plethysm multiplicity, the connected-stabilizer rectangular Kronecker multiplicity, and the coordinate-ring multiplicity of an orbit closure. The transpose component of the determinant stabilizer and the orbit boundary are separate restrictions. R-03 gives finite counterexamples to identifying them. R-05 computes the closure multiplicity by an actual rank, and R-08 expresses the same requirement as a quotient numerator.

If $X\subseteq W$ is $GL(V)$-stable and $V_k$ is spanned by the first $k$ basis vectors, let $X_k$ be the closure of its restriction image in $\operatorname{Sym}^nV_k^*$. For $\ell(\lambda)\leq k$, the multiplicities in $\mathbb C[X]$ and $\mathbb C[X_k]$ agree. The short verification used repeatedly below is as follows. A highest-weight polynomial of weight $(\lambda_1,\ldots,\lambda_k,0,\ldots,0)$ uses only coefficient variables supported on $V_k$, since all generator weights are nonnegative. It vanishes on $X$ exactly when it vanishes on $X_k$. Taking the quotient of the common ambient highest-weight space proves the assertion. This is a coefficient-subalgebra version of the usual inheritance principle, proved here in precisely the form needed for the orbit restrictions.

The exact identities alone impose no efficient complexity bound. In particular, a nonnegative squared-norm formula may have exponentially many terms; a rank formula may require exponentially large matrices; and a positive dimension may become zero after highest-weight projection. The size and sign statements in each result are therefore part of the theorem, not suppressed implementation details.

\subsection{How the specific new analytic tools enter}

Alp\"oge's step~2 (p.~3) concerns theta derivatives and inert-prime divisibility. It does not, in the supplied four-page text, state a bound for a coefficient-matroid minor. Steps~3 and~5 (p.~3) motivate trace and square-root-cancellation arguments; P-02, P-04 and P-06 implement the corresponding curve/sheaf constructions with their exact degree and conductor budgets. Steps~4 and~7 (pp.~3--4) concern determinant and Gan--Gross--Prasad period positivity for specified groups; R-03 identifies why that branching problem cannot be substituted for the tensor subgroup without a proved functorial bridge. Step~8 (p.~4) concerns theta-lift multiplicities; I-03 and I-14 identify the different coefficient-span quantifiers a proposed isolation transfer must preserve.

Step~13 (p.~4) invokes a Bost slope inequality, and step~17 (p.~4) concerns the displayed arithmetic cohomology groups. The supplied note gives no reusable estimates for the ranks, heights or valuations of the ROABP coefficient maps used here. Their possible use therefore has a concrete prerequisite: construct the relevant arithmetic bundle or cohomology object, identify the desired minor or quotient with its invariant, and retain effective size bounds. This is a dependency statement about the supplied inputs, not a claim that such a construction is impossible. The integral and special-fiber analyses below identify where that next construction would act.
